% Schnittmengen – bank/schnittmengen/ (zusammenbau v0.4, Heft)
\einheitenkopf[t6]{Schnittmengen}
\setcounter{aufgabe}{56}

% Anlauf (ohne Prüfkennung)
\begin{aufgabe}{Schnittpunkt von Gerade und Ebene – Anlauf}
\begin{teile}
\teil Gegeben sind eine Gerade in Parameterform und eine Ebene in Koordinatenform. Welcher Weg führt zum Schnittpunkt? \\ \kreuz{einsetzen} \\ \kreuz{gleichsetzen}
\teil Bestimme den Schnittpunkt der Geraden $\vec{x} = (-1|2|1) + t \cdot (1|0|1)$ mit der Ebene $x + y + z = 6$. S( \leerfeld | \leerfeld | \leerfeld )
\end{teile}
\end{aufgabe}

\begin{aufgabe}{Schnittpunkt von Gerade und Ebene – Prüfungsaufgaben}
\begin{teile}
\teil Ein Körper hat die Spitze $P(0|0|3)$; $M(3|1|0)$ ist der Mittelpunkt einer Grundkante. Die Gerade durch $P$ und $M$ schneidet die Ebene $L$: $x + 2y + z = 4$ im Punkt $Q$. Bestimme die Koordinaten von $Q$. \hfill (Abitur 2023 GK) \\ Q( \leerfeld | \leerfeld | \leerfeld )
\teil Aus einem Versorgungsstollen führt von $L(4|6|2)$ ein Rohr senkrecht nach oben durch den Berg und ragt 1,5 m aus ihm heraus. Die Bergoberfläche liegt in der Ebene $F$: $2x + 3y + 5z = 51$; die $xy$-Ebene liegt auf Meereshöhe, eine Längeneinheit entspricht 50 m. Berechne die Länge des Rohrs. \hfill (Abitur 2019 GK) \\ \leerfeld[m]
\teil Ein Sonnensegel hat die Ecken $A(1|2|2)$, $B(1|6|2)$ und $C(-2|4|4)$; der Boden ist die $xy$-Ebene, eine Längeneinheit entspricht 1 m. Bei parallelem Sonnenlicht liegen die Schatten von $A$ und $B$ in $A'(-1|3|0)$ und $B'(-1|7|0)$. Bestimme den Schatten $C'$ von $C$. \hfill (Abitur 2020 GK) \\ C'( \leerfeld | \leerfeld | \leerfeld )
\teil Durch die Glasfassade eines Hauses fällt Sonnenlicht in Richtung $(2|-1|-2)$; eine Längeneinheit entspricht 1 m. Der Boden des Dachgeschosses liegt in der Ebene $z = 6$, der obere Eckpunkt der Fassade ist $L(0|5|9)$. Bestimme den Punkt, in dem das Licht durch $L$ auf den Boden des Dachgeschosses trifft. \hfill (Abitur 2019 GK) \\ ( \leerfeld | \leerfeld | \leerfeld )
\teil Neben einem Fahnenmast mit der Spitze $S(-1|1|5)$ steht ein Pfosten; sein oberes Ende $P(2|3|2)$ wirft bei parallelem Sonnenlicht den Schatten $P'(3|2|0)$ auf den Boden ($xy$-Ebene). Bestimme den Schatten $S'$ der Mastspitze. \hfill (Abitur 2020 GK) \\ S'( \leerfeld | \leerfeld | \leerfeld )
\teil Ein Rasenroboter fährt geradlinig und gleichförmig von $P(8|2|0)$ über $Q(6|3|0)$ weiter; für die Strecke von $P$ nach $Q$ braucht er 4 s. Eine Lichtschranke liegt in der senkrechten Ebene $E$: $x - 2y = -6$. Wie lange braucht der Roboter von $Q$ bis zur Lichtschranke? \hfill (Abitur 2018 GK) \\ \leerfeld[s]
\teil Ein Boot fährt geradlinig mit konstanter Geschwindigkeit von $P(1|8|0)$ nach $Q(4|6|0)$ und braucht dafür 2 Minuten; es fährt geradeaus weiter. Die Grenze eines Schutzgebiets liegt in der senkrechten Ebene $E$: $2x + y = 20$. Wie lange braucht das Boot von $Q$ bis zur Grenze? \hfill (Abitur 2018 GK) \\ \leerfeld[min]
\teil Eine Drohne fliegt geradlinig und gleichförmig von $P(2|1|4)$ über $Q(3|3|5)$ weiter; von $P$ nach $Q$ braucht sie 5 s. Eine Hauswand liegt in der Ebene $E$: $x + y = 15$. Wie lange braucht die Drohne von $Q$ bis zur Hauswand? \hfill (Abitur 2018 GK) \\ \leerfeld[s]
\end{teile}
\end{aufgabe}

% Anlauf (ohne Prüfkennung)
\begin{aufgabe}{Schnittpunkt zweier Geraden – Anlauf}
\begin{teile}
\teil Zwei Geraden im Raum werden gleichgesetzt: $(1|0|2) + r \cdot (1|2|0) = (3|4|1) + s \cdot (0|1|1)$. Was entsteht? \\ \kreuz{drei Gleichungen mit zwei Unbekannten} \\ \kreuz{zwei Gleichungen mit zwei Unbekannten} \\ \kreuz{drei Gleichungen mit drei Unbekannten}
\teil Bestimme den Schnittpunkt der Geraden $g: \vec{x} = (2|2|3) + r \cdot (1|0|2)$ und $h: \vec{x} = (3|0|3) + s \cdot (0|1|1)$. S( \leerfeld | \leerfeld | \leerfeld )
\end{teile}
\end{aufgabe}

\begin{aufgabe}{Schnittpunkt zweier Geraden – Prüfungsaufgaben}
\begin{teile}
\teil Ein Mähroboter startet auf einer ebenen, leicht geneigten Rasenfläche im Punkt $P(2|4{,}5|0{,}25)$ und fährt entlang der Geraden durch $P$ mit dem Richtungsvektor $(4|-1|0{,}5)$ auf den Rand $BC$ zu; $B(8|0|1)$, $C(8|6|1)$, eine Längeneinheit entspricht 1 m. Bestimme den Punkt $Q$, in dem die Bahn die Strecke $BC$ trifft (zur Kontrolle: $Q(8|3|1)$). \hfill (Abitur 2021 GK) \\ Q( \leerfeld | \leerfeld | \leerfeld )
\teil Gegeben sind $g: \vec{x} = (1|4|-2) + s \cdot (2|0|3)$ und die Gerade $h$ durch $A(3|0|1)$ und $B(4|1|b)$ mit einer reellen Zahl $b$. $g$ und $h$ haben einen gemeinsamen Punkt. Bestimme $b$. \hfill (Abitur 2023 GK) \\ b = \leerfeld
\teil Gegeben sind $g: \vec{x} = (0|-2|5) + s \cdot (1|0|-1)$ und die Gerade $h$ durch $A(2|1|0)$ und $B(3|2|b)$ mit einer reellen Zahl $b$. $g$ und $h$ haben einen gemeinsamen Punkt. Bestimme $b$. \hfill (Abitur 2023 GK) \\ b = \leerfeld
\teil Gegeben sind $g: \vec{x} = (3|1|0) + s \cdot (0|2|3)$ und die Gerade $h$ durch $A(1|3|2)$ und $B(2|3|b)$ mit einer reellen Zahl $b$. $g$ und $h$ haben einen gemeinsamen Punkt. Bestimme $b$. \hfill (Abitur 2023 GK) \\ b = \leerfeld
\teil Gegeben ist der Würfel $ABCDEFGH$ mit $A(0|0|0)$, $G(4|4|4)$ und $H(0|4|4)$. Für jede reelle Zahl $r$ ist die Gerade $g_r: \vec{x} = (1 - r|0|r^2) + u \cdot (2|4|0)$, $u \in \mathbb{R}$, gegeben. Für einen Wert von $r$ schneidet $g_r$ die Kante $GH$. Bestimme das Verhältnis, in dem der Schnittpunkt die Kante $GH$ teilt. \hfill (Abitur 2019 GK) \\ HQ : QG = \leerfeld : \leerfeld
\teil Gegeben ist der Würfel $ABCDEFGH$ mit $A(0|0|0)$, $G(6|6|6)$ und $H(0|6|6)$. Für jede reelle Zahl $r$ ist die Gerade $g_r: \vec{x} = (3 + r|0|r^2 - 3) + u \cdot (1|3|0)$, $u \in \mathbb{R}$, gegeben. Für einen Wert von $r$ schneidet $g_r$ die Kante $GH$. Bestimme das Verhältnis, in dem der Schnittpunkt die Kante $GH$ teilt. \hfill (Abitur 2019 GK) \\ HQ : QG = \leerfeld : \leerfeld
\teil Gegeben ist der Würfel $ABCDEFGH$ mit $A(0|0|0)$, $G(2|2|2)$ und $H(0|2|2)$. Für jede reelle Zahl $r$ ist die Gerade $g_r: \vec{x} = (1 - 0{,}5r|0|r^2 - 2) + u \cdot (1|2|0)$, $u \in \mathbb{R}$, gegeben. Für einen Wert von $r$ schneidet $g_r$ die Kante $GH$. Bestimme das Verhältnis, in dem der Schnittpunkt die Kante $GH$ teilt. \hfill (Abitur 2019 GK) \\ HQ : QG = \leerfeld : \leerfeld
\teil Gegeben ist der Würfel $ABCDEFGH$ mit $A(0|0|0)$, $G(8|8|8)$ und $H(0|8|8)$. Für jede reelle Zahl $r$ ist die Gerade $g_r: \vec{x} = (7 + 2r|0|2r^2) + u \cdot (1|4|0)$, $u \in \mathbb{R}$, gegeben. Für einen Wert von $r$ schneidet $g_r$ die Kante $GH$. Bestimme das Verhältnis, in dem der Schnittpunkt die Kante $GH$ teilt. \hfill (Abitur 2019 GK) \\ HQ : QG = \leerfeld : \leerfeld
\end{teile}
\end{aufgabe}

\begin{aufgabe}{Schattenpunkt auf einer Kante als Schnittpunkt von Lichtgerade und Kantengerade berechnen – Prüfungsaufgaben}
\begin{teile}
\teil Ein Spielturm hat ein Dach in Form einer Pyramide; an der Spitze ist eine senkrechte Stange mit dem oberen Ende $T(0|0|7)$ befestigt; der Boden ist die $xy$-Ebene, eine Längeneinheit entspricht 1 m. Sonnenlicht fällt in parallelen Geraden mit dem Richtungsvektor $(4|-1|-2)$. Der Schatten des oberen Endes fällt auf die Dachkante $EF$ mit $E(6|-3|4)$ und $F(6|3|4)$. Bestimme diesen Schattenpunkt. \hfill (Abitur 2017 GK) \\ ( \leerfeld | \leerfeld | \leerfeld )
\teil An einem Pavillon ist eine Wetterfahne mit dem oberen Ende $T(1|1|7)$ befestigt; der Boden ist die $xy$-Ebene, eine Längeneinheit entspricht 1 m. Sonnenlicht fällt in parallelen Geraden mit dem Richtungsvektor $(2|1|-3)$. Der Schatten des oberen Endes fällt auf die Dachkante $EF$ mit $E(3|-1|4)$ und $F(3|5|4)$. Bestimme diesen Schattenpunkt. \hfill (Abitur 2017 GK) \\ ( \leerfeld | \leerfeld | \leerfeld )
\teil Auf einem Gartenhaus steht ein Blitzableiter mit der Spitze $T(0|0|4{,}5)$; der Boden ist die $xy$-Ebene, eine Längeneinheit entspricht 1 m. Sonnenlicht fällt in parallelen Geraden mit dem Richtungsvektor $(-2|2|-1)$. Der Schatten des oberen Endes fällt auf die Dachkante $EF$ mit $E(-3|1|4)$ und $F(1|1|4)$. Bestimme diesen Schattenpunkt. \hfill (Abitur 2017 GK) \\ ( \leerfeld | \leerfeld | \leerfeld )
\end{teile}
\end{aufgabe}

% Anlauf (ohne Prüfkennung)
\begin{aufgabe}{Spuren, Schnittgeraden und Schnittfiguren – Anlauf}
\begin{teile}
\teil Der Spurpunkt einer Ebene auf der $y$-Achse – welche Koordinaten sind null? \\ \kreuz{$x$ und $z$} \\ \kreuz{nur $y$} \\ \kreuz{nur $x$}
\teil Bestimme die Spurpunkte der Ebene $2x + 3y + 6z = 6$ auf den Koordinatenachsen. $S_x($ \leerfeld |0|0), $S_y(0|$ \leerfeld |0), $S_z(0|0|$ \leerfeld )
\end{teile}
\end{aufgabe}

\begin{aufgabe}{Spuren, Schnittgeraden und Schnittfiguren – Prüfungsaufgaben}
\begin{teile}
\teil Bestimme die Spurpunkte der Ebene $6x + 4y + 3z = 24$ auf den Koordinatenachsen. \hfill (Abitur 2020 GK) \\ $S_x($ \leerfeld |0|0), $S_y(0|$ \leerfeld |0), $S_z(0|0|$ \leerfeld )
\teil Ein Holzkörper ist eine Pyramide über dem Quadrat $ABCD$ mit $A(0|0|0)$, $B(5|0|0)$, $C(5|5|0)$, $D(0|5|0)$ und der Spitze $E(0|5|4)$ senkrecht über $D$; eine Längeneinheit entspricht 1 cm. Die Ebene $L$: $4x + 5z = 20$ schneidet die $z$-Achse im Punkt $F$. Bestimme $F$ und zeichne die Schnittgeraden von $L$ mit der $xz$-Ebene und mit der $yz$-Ebene in das Schrägbild ein. \hfill (Abitur 2021 GK)

\begin{ksys3}[xyz,x1min=0,x1max=6,x2min=0,x2max=6,x3min=0,x3max=5] \rpyramide{0,0,0}{5}{5}{0,5,4} \end{ksys3}
\teil Ein Holzkörper ist eine Pyramide über dem Quadrat $ABCD$ mit $A(0|0|0)$, $B(7|0|0)$, $C(7|7|0)$, $D(0|7|0)$ und der Spitze $E(0|7|3)$ senkrecht über $D$; eine Längeneinheit entspricht 1 cm. Die Ebene $L$: $3x + 7z = 21$ schneidet die $z$-Achse im Punkt $F$. Bestimme $F$ und zeichne die Schnittgeraden von $L$ mit der $xz$-Ebene und mit der $yz$-Ebene in das Schrägbild ein. \hfill (Abitur 2021 GK)

\begin{ksys3}[xyz,x1min=0,x1max=8,x2min=0,x2max=8,x3min=0,x3max=4] \rpyramide{0,0,0}{7}{7}{0,7,3} \end{ksys3}
\teil Die Abbildung zeigt den Quader mit $A(7|0|0)$, $B(7|7|0)$, $C(0|7|0)$ und $F(7|7|2)$, die Gerade $h$ durch $B$ und $F$ sowie die Schnittfigur des Quaders mit einer Ebene durch $A$, $C$ und einen Punkt $P$ von $h$. Beschreibe, wie man mithilfe der Abbildung ermitteln kann, dass $P$ die $z$-Koordinate 4 hat. \hfill (Abitur 2018 GK)

\begin{ksys3}[xyz,x1min=0,x1max=8,x2min=0,x2max=8,x3min=0,x3max=5] \rquader{0,0,0}{7}{7}{2} \rgerade[0:1]{7,0,0}{0,3.5,2}{} \rgerade[0:1]{7,3.5,2}{-3.5,3.5,0}{} \rgerade[0:1]{3.5,7,2}{-3.5,0,-2}{} \rgerade[0:1]{0,7,0}{7,-7,0}{} \rgerade[0:5]{7,7,0}{0,0,1}{h} \end{ksys3}
\teil Die Abbildung zeigt den Quader mit $A(5|0|0)$, $B(5|5|0)$, $C(0|5|0)$ und $F(5|5|2)$, die Gerade $h$ durch $B$ und $F$ sowie die Schnittfigur des Quaders mit einer Ebene durch $A$, $C$ und einen Punkt $P$ von $h$. Beschreibe, wie man mithilfe der Abbildung ermitteln kann, dass $P$ die $z$-Koordinate 4 hat. \hfill (Abitur 2018 GK)

\begin{ksys3}[xyz,x1min=0,x1max=6,x2min=0,x2max=6,x3min=0,x3max=5] \rquader{0,0,0}{5}{5}{2} \rgerade[0:1]{5,0,0}{0,2.5,2}{} \rgerade[0:1]{5,2.5,2}{-2.5,2.5,0}{} \rgerade[0:1]{2.5,5,2}{-2.5,0,-2}{} \rgerade[0:1]{0,5,0}{5,-5,0}{} \rgerade[0:5]{5,5,0}{0,0,1}{h} \end{ksys3}
\end{teile}
\end{aufgabe}

\begin{aufgabe}{Spuren, Schnittgeraden und Schnittfiguren – Prüfungsaufgaben – weiter}
\begin{teile}
\teil Gegeben ist das Quadrat $ABCD$ mit $A(0|0|0)$, $B(8|0|0)$, $C(8|8|0)$ und $D(0|8|0)$ sowie die Ebenen $E_1$: $x = 4$, $E_2$: $x + y = 8$ und $E_3$: $y = 2$. Entscheide für jede der drei Ebenen, ob sie das Quadrat in zwei Figuren mit gleichem Flächeninhalt teilt, und begründe. \hfill (Abitur 2021 GK)
\teil Gegeben ist das Quadrat $ABCD$ mit $A(0|0|0)$, $B(5|0|0)$, $C(5|5|0)$ und $D(0|5|0)$ sowie die Ebenen $E_1$: $z = 1$, $E_2$: $x - y = 0$ und $E_3$: $x = 2$. Entscheide für jede der drei Ebenen, ob sie das Quadrat in zwei Figuren mit gleichem Flächeninhalt teilt, und begründe. \hfill (Abitur 2021 GK)
\teil Gegeben ist das Quadrat $ABCD$ mit $A(1|1|0)$, $B(7|1|0)$, $C(7|7|0)$ und $D(1|7|0)$ sowie die Ebenen $E_1$: $y = 4$, $E_2$: $x + y = 8$ und $E_3$: $x = 5$. Entscheide für jede der drei Ebenen, ob sie das Quadrat in zwei Figuren mit gleichem Flächeninhalt teilt, und begründe. \hfill (Abitur 2021 GK)
\teil Gegeben sind die Ebenen $E_1: \vec{x} = (5|0|0) + r \cdot (-1|1|0) + s \cdot (-1|0|1)$ und $E_2$: $2x + 2y + z = 10$. Begründe, dass $E_1$ und $E_2$ nicht parallel sind, und bestimme eine Gleichung ihrer Schnittgeraden. \hfill (Abitur 2020 GK) \\ x = \leerfeld + r · \leerfeld
\teil Gegeben sind die Ebenen $E_1: \vec{x} = (0|3|0) + r \cdot (2|-3|0) + s \cdot (0|-3|1)$ und $E_2$: $2x + y + 4z = 4$. Begründe, dass $E_1$ und $E_2$ nicht parallel sind, und bestimme eine Gleichung ihrer Schnittgeraden. \hfill (Abitur 2020 GK) \\ x = \leerfeld + r · \leerfeld
\end{teile}
\end{aufgabe}
